 \section{非负矩阵分解}

许多统计数据都可以自然地组织成一个非负矩阵。例如，在文本分析中，矩阵的行可以表示词语，列可以表示文档，矩阵元素表示某个词在某篇文档中的出现频数；在图像分析中，矩阵元素可以表示像素强度；在推荐系统中，矩阵元素可以表示用户对商品的评分或交互次数；在基因表达分析中，矩阵元素可以表示某个基因在某个样本或细胞中的表达强度。虽然这些数据来自不同问题，但它们有一个共同特点：观测矩阵往往维度很高，而真正支配数据变化的结构通常远少于原始维度。\par
因此，一个自然的统计问题是：能否用少数潜在成分来近似解释原始数据？设观测数据矩阵$X=(x_{ij})\in\mathbb{R}^{n\times p}$，且$x_{ij}\geqslant0$，\gls{NMF}希望寻找两个非负矩阵：
\begin{equation*}
	W=(w_{ik})\in\mathbb{R}^{n\times r},\quad H=(h_{kj})\in\mathbb{R}^{r\times p},\quad w_{ik}\geqslant0,\quad h_{kj}\geqslant0
\end{equation*}
使得
\begin{equation*}
	X\approx WH,\quad r<\min\{n,p\}
\end{equation*}
这里的分解并不只是为了压缩数据，更重要的是希望把复杂观测拆解为少数可解释的潜在结构。$r$称为分解秩或潜在因子数。每个样本列向量$x_j$都有近似表示：
\begin{equation*}
	x_j\approx Wh_j=\sum_{k=1}^{r}h_{kj}w_k,\quad h_{kj}\geqslant0
\end{equation*}
直观上，$W$的第$k$列可理解为第$k$个非负基向量，$H$的第$k$行给出该基向量在各个样本中的非负系数。换言之，矩阵分解试图回答的问题是：原始高维数据是否可以由少数共同模式组合而成？\par
经典的低秩分解方法，如PCA或SVD，也是在回答类似问题。它们通过寻找低维线性子空间来获得最优的线性低秩近似，从而在平方误差意义下有效保留数据中的主要变化方向。然而，对于天然非负的数据，这类方法存在一个解释上的困难：分解得到的基向量和系数通常允许取负值。负值在代数上没有问题，甚至对最优线性近似是有利的，因为它允许不同方向之间通过正负抵消来表示数据；但在许多实际问题中，这种抵消机制并不符合数据生成过程。例如，词频不能解释为某些主题词的正贡献再减去另一些主题词的负贡献，像素强度也不能自然地解释为若干局部图像部件之间的正负抵消，基因表达强度同样更适合被理解为若干表达程序的非负叠加。\par
\begin{definition}
	设$X=(x_{ij})\in\mathbb{R}^{n\times p}$且$x_{ij}\geqslant0$，给定正整数$r<\min\{m,n\}$。若存在$W\in\mathbb{R}^{n\times r},\;H\in\mathbb{R}^{p\times n}$满足：
	\begin{equation*}
		X\approx WH,\quad W=(w_{ik})\in\mathbb{R}^{n\times r},\quad H=(h_{kj})\in\mathbb{R}^{r\times p},\quad w_{ik}\geqslant0,\quad h_{kj}\geqslant0
	\end{equation*}
	则称$X\approx WH$为$X$的非负矩阵分解。其中$W$称为基矩阵，$H$称为系数矩阵。
\end{definition}
\begin{property}\label{prop:NMF}
	非负矩阵分解具有如下基本性质：
	\begin{enumerate}
		\item 分解不唯一。
		\item $\operatorname{rank}(WH)\leqslant r$，即非负矩阵分解是一种带有非负约束的低秩近似方法。
	\end{enumerate}
\end{property}
\begin{proof}
	(1)改变分解矩阵行列的排列顺序或对分解矩阵进行可逆变换即可给出同一矩阵的多种分解结果。\par
	(4)由\cref{prop:Rank}(9)立即可得。
\end{proof}

\subsection{目标函数}
\subsubsection{Frobenius范数目标}
当数据近似服从加性高斯噪声模型时，最常用的目标函数是平方误差：
\begin{equation*}
	\min_{W,H}\;F(W,H)=\frac{1}{2}||X-WH||_F^2,\quad W\geqslant0,\quad H\geqslant0
\end{equation*}
其中：
\begin{equation*}
	||A||_F^2=\sum_{i=1}^{m}\sum_{j=1}^{n}a_{ij}^2=\operatorname{tr}(A^{\top}A)
\end{equation*}
称为矩阵$A$的Frobenius范数平方。

\begin{derivation}
	对$F(W,H)$求导可得：
	\begin{align*}
		\nabla_WF(W,H)&=(WH-X)H^{\top}=WHH^{\top}-XH^{\top} \\
		\nabla_HF(W,H)&=W^{\top}(WH-X)=W^{\top}WH-W^{\top}X
	\end{align*}
	若暂时不考虑非负约束，则固定$W$时关于$H$的最优解满足正规方程：
	\begin{equation*}
		W^{\top}WH=W^{\top}X
	\end{equation*}
	固定$H$时关于$W$的最优解满足：
	\begin{equation*}
		WHH^{\top}=XH^{\top}
	\end{equation*}
	非负约束使得这两个线性方程不能直接给出最终解，但它们解释了许多迭代算法的来源：每次固定一边，更新另一边，使得重构误差逐步下降。
\end{derivation}

\subsubsection{KL散度目标}
当$X$是计数矩阵或频数矩阵时，常用如下广义KL散度度量$X$和$Y=WH$之间的差异：
\begin{equation*}
	D(X\|WH)=\sum_{i=1}^{m}\sum_{j=1}^{n}\left[x_{ij}\ln\frac{x_{ij}}{(WH)_{ij}}-x_{ij}+(WH)_{ij}\right]
\end{equation*}
其中约定$0\ln0=0$。这一目标与Poisson模型关系密切：若$x_{ij}$近似服从均值为$(WH)_{ij}$的Poisson分布且各元素独立，则最大化似然等价于最小化$D(X\|WH)$。

\begin{note}
	Frobenius目标更接近连续型强度数据的平方误差近似，KL散度目标更适合非负计数数据或相对比例结构明显的数据。实际使用时，目标函数的选择不只是计算问题，也是在选择误差模型。
\end{note}

\subsection{负二项分布NMF及其gene dependence扩展}
当$X$为基因--样本计数矩阵时，Poisson模型所要求的均值与方差相等通常过于严格。负二项分布允许方差随均值更快增长，因此可以在保留非负低秩均值结构的同时描述过度离散的计数噪声。沿用前文记号，以下将第$i$行视为第$i$个gene，将第$j$列视为第$j$个样本，且$x_{ij}\in\mathbb{N}_0$。

\subsubsection{Mean-dispersion参数化}
\begin{definition}
	设$X$为非负整数值随机变量，$\mu>0,\;\phi>0$。若$X$服从mean-dispersion参数化的负二项分布，则记作$X\sim\operatorname{NB}(\mu,\phi)$，其中：
	\begin{equation*}
		\operatorname{E}(X)=\mu,\qquad \operatorname{Var}(X)=\mu+\phi\mu^2.
	\end{equation*}
	$\mu$称为均值参数，$\phi$称为离散参数。
\end{definition}

\begin{derivation}
	从参数为$r>0,\;p\in(0,1)$的负二项分布出发，其概率质量函数为：
	\begin{equation*}
		P(X=x)=\frac{\Gamma(x+r)}{\Gamma(r)\Gamma(x+1)}p^r(1-p)^x,\quad x\in\mathbb{N}_0.
	\end{equation*}
	令：
	\begin{equation*}
		r=\phi^{-1},\qquad p=\frac{r}{r+\mu}=\frac{1}{1+\phi\mu},\qquad 1-p=\frac{\phi\mu}{1+\phi\mu}.
	\end{equation*}
	由$\operatorname{E}(X)=r(1-p)/p$和$\operatorname{Var}(X)=r(1-p)/p^2$可得：
	\begin{align*}
		\operatorname{E}(X)&=\mu, \\
		\operatorname{Var}(X)&=\mu(1+\phi\mu)=\mu+\phi\mu^2.
	\end{align*}
	因此，mean-dispersion参数化下的概率质量函数为：
	\begin{equation*}
		f_{\operatorname{NB}}(x;\mu,\phi)
		=\frac{\Gamma(x+\phi^{-1})}{\Gamma(\phi^{-1})\Gamma(x+1)}
		\left(\frac{1}{1+\phi\mu}\right)^{\phi^{-1}}
		\left(\frac{\phi\mu}{1+\phi\mu}\right)^x.
	\end{equation*}
	单个观测值的对数似然为：
	\begin{align*}
		\ell(x;\mu,\phi)
		={}&\log\Gamma(x+\phi^{-1})-\log\Gamma(\phi^{-1})-\log\Gamma(x+1) \\
		&+x\log(\phi\mu)-(x+\phi^{-1})\log(1+\phi\mu).
	\end{align*}
\end{derivation}

\subsubsection{原始NB-NMF}
沿用$X\approx WH$的分解形式，令负二项分布的均值结构为：
\begin{equation*}
	\mu_{ij}=(WH)_{ij}=\sum_{k=1}^{r}w_{ik}h_{kj},\quad i=1,2,\dots,n,\quad j=1,2,\dots,p.
\end{equation*}
为允许不同gene具有不同的过度离散程度，记$\phi=(\phi_1,\phi_2,\dots,\phi_n)^{\top}\in(0,+\infty)^n$，并假设：
\begin{equation*}
	x_{ij}\mid W,H,\phi\overset{\mathrm{ind}}{\sim}\operatorname{NB}(\mu_{ij},\phi_i).
\end{equation*}
若进一步假设所有gene具有相同的离散程度，只需令$\phi_i=\phi$。在上述逐元素条件独立假设下，原始NB-NMF的似然和对数似然分别为：
\begin{align*}
	L_{\operatorname{NB}}(W,H,\phi;X)
	&=\prod_{i=1}^{n}\prod_{j=1}^{p}f_{\operatorname{NB}}\left(x_{ij};(WH)_{ij},\phi_i\right), \\
	\ell_{\operatorname{NB}}(W,H,\phi;X)
	&=\sum_{i=1}^{n}\sum_{j=1}^{p}\ell\left(x_{ij};(WH)_{ij},\phi_i\right).
\end{align*}
记参数约束集合为：
\begin{equation*}
	\mathcal{C}_{\operatorname{NB}}
	=\left\{(W,H,\phi):
	\begin{aligned}
		&W\in\mathbb{R}^{n\times r},\;H\in\mathbb{R}^{r\times p},\;W\geqslant0,\;H\geqslant0,\\
		&\phi\in(0,+\infty)^n,\;(WH)_{ij}>0,\;i=1,2,\dots,n,\;j=1,2,\dots,p
	\end{aligned}
	\right\}.
\end{equation*}
则原始NB-NMF的极大似然估计问题为：
\begin{equation*}
	(\widehat{W},\widehat{H},\widehat{\phi})
	\in\operatorname*{argmax}_{(W,H,\phi)\in\mathcal{C}_{\operatorname{NB}}}
	\ell_{\operatorname{NB}}(W,H,\phi;X).
\end{equation*}

\begin{note}
	逐元素条件独立是一个working likelihood假设，而不是断言不同gene在生物学上真实独立。$W$的列表示共同gene program，$H$的列表示各样本的program强度，因此低秩均值$WH$本身已经能够解释一部分gene--gene共表达。原始NB-NMF忽略的是：在给定$WH$之后，尚未被低秩均值结构吸收的residual gene dependence。
\end{note}

\begin{derivation}
	固定$\phi$，单个对数似然项关于均值参数$\mu$的score为：
	\begin{align*}
		\frac{\partial\ell(x;\mu,\phi)}{\partial\mu}
		&=\frac{x}{\mu}-\frac{(x+\phi^{-1})\phi}{1+\phi\mu} \\
		&=\frac{x-\mu}{\mu(1+\phi\mu)}
		=\frac{x-\mu}{\mu+\phi\mu^2}.
	\end{align*}
	记$S=(s_{ij})\in\mathbb{R}^{n\times p}$，其中：
	\begin{equation*}
		s_{ij}=\left.\frac{\partial\ell(x_{ij};\mu,\phi_i)}{\partial\mu}\right|_{\mu=(WH)_{ij}}
		=\frac{x_{ij}-(WH)_{ij}}{(WH)_{ij}+\phi_i(WH)_{ij}^2}.
	\end{equation*}
	由$\partial\mu_{ij}/\partial w_{ik}=h_{kj}$和$\partial\mu_{ij}/\partial h_{kj}=w_{ik}$，链式法则给出：
	\begin{align*}
		\frac{\partial\ell_{\operatorname{NB}}}{\partial w_{ik}}
		&=\sum_{j=1}^{p}s_{ij}h_{kj}, &
		\frac{\partial\ell_{\operatorname{NB}}}{\partial h_{kj}}
		&=\sum_{i=1}^{n}w_{ik}s_{ij}, \\
		\nabla_W\ell_{\operatorname{NB}}
		&=SH^{\top}\in\mathbb{R}^{n\times r}, &
		\nabla_H\ell_{\operatorname{NB}}
		&=W^{\top}S\in\mathbb{R}^{r\times p}.
	\end{align*}
	这里的梯度均在固定$\phi$时计算。
\end{derivation}

非负约束使score方程不能直接给出闭式解。此外，对任意正对角矩阵$C\in\mathbb{R}^{r\times r}$，$WH=(WC)(C^{-1}H)$，所以$W$与$H$之间还存在scale non-identifiability。实际计算通常需要结合投影梯度、坐标更新或其它带约束的数值优化方法，并通过列归一化等约束固定尺度。

\subsubsection{Gene graph regularized NB-NMF}
若gene之间存在已知网络结构，可以在原始NB-NMF的估计中加入graph Laplacian penalty。设$A=(a_{ii'})\in\mathbb{R}^{n\times n}$为给定的无向gene graph邻接矩阵，其中$A=A^{\top},\;a_{ii'}\geqslant0$。为了与前文的KL散度$D(X\|WH)$区分，记度矩阵为$\Delta\in\mathbb{R}^{n\times n}$，并定义：
\begin{equation*}
	\Delta_{ii}=\sum_{i'=1}^{n}a_{ii'},\qquad \Delta_{ii'}=0\;(i\neq i'),\qquad L=\Delta-A.
\end{equation*}
由于gene位于$X$和$W$的行方向，第$i$个gene在program space中的loading是$W$的第$i$行$w_{i\cdot}\in\mathbb{R}^{r}$，所以图惩罚应作用于$W$：
\begin{equation*}
	\mathcal{R}_{\operatorname{graph}}(W)=\operatorname{tr}(W^{\top}LW).
\end{equation*}
其中$W^{\top}LW\in\mathbb{R}^{r\times r}$，故该迹运算的维度相容。

\begin{derivation}
	由$L=\Delta-A$以及$A=A^{\top}$可得：
	\begin{align*}
		\operatorname{tr}(W^{\top}LW)
		&=\sum_{i=1}^{n}\Delta_{ii}||w_{i\cdot}||_2^2
		-\sum_{i=1}^{n}\sum_{i'=1}^{n}a_{ii'}w_{i\cdot}w_{i'\cdot}^{\top} \\
		&=\frac{1}{2}\sum_{i=1}^{n}\sum_{i'=1}^{n}a_{ii'}
		||w_{i\cdot}-w_{i'\cdot}||_2^2.
	\end{align*}
	因此，若两个gene在图上具有较大的连接权重，则较大的loading差异会受到更强惩罚。
\end{derivation}

给定正则化参数$\lambda\geqslant0$，定义penalized log-likelihood：
\begin{equation*}
	\ell_{\operatorname{graph}}(W,H,\phi;X)
	=\ell_{\operatorname{NB}}(W,H,\phi;X)
	-\lambda\operatorname{tr}(W^{\top}LW).
\end{equation*}
由于图惩罚只作用于$W$，若仍允许$W\mapsto cW,\;H\mapsto H/c$，则在保持$WH$不变的同时可以任意缩小惩罚。为固定这一尺度，记：
\begin{equation*}
	\mathcal{C}_{\operatorname{graph}}
	=\left\{(W,H,\phi)\in\mathcal{C}_{\operatorname{NB}}:
	\sum_{i=1}^{n}w_{ik}=1,\;k=1,2,\dots,r\right\}.
\end{equation*}
相应的penalized maximum likelihood estimation问题为：
\begin{equation*}
	(\widehat{W}_{\operatorname{graph}},\widehat{H}_{\operatorname{graph}},\widehat{\phi}_{\operatorname{graph}})
	\in\operatorname*{argmax}_{(W,H,\phi)\in\mathcal{C}_{\operatorname{graph}}}
	\ell_{\operatorname{graph}}(W,H,\phi;X).
\end{equation*}

\begin{derivation}
	对一般矩阵$L$，有：
	\begin{equation*}
		\nabla_W\operatorname{tr}(W^{\top}LW)=(L+L^{\top})W.
	\end{equation*}
	无向图的Laplacian矩阵满足$L=L^{\top}$，故：
	\begin{equation*}
		\nabla_W\operatorname{tr}(W^{\top}LW)=2LW.
	\end{equation*}
	因此，图正则化仅在$W$的梯度中加入修正项：
	\begin{align*}
		\nabla_W\ell_{\operatorname{graph}}&=SH^{\top}-2\lambda LW, \\
		\nabla_H\ell_{\operatorname{graph}}&=W^{\top}S.
	\end{align*}
\end{derivation}

\begin{note}
	Gene graph regularized NB-NMF并没有构造一个完整的多元负二项分布，其似然部分仍是逐元素负二项概率质量函数的乘积。该模型通过graph Laplacian penalty把外部gene graph作为先验结构注入program loading的估计中。它的形式简单、解释直接，计算复杂度也接近原始NB-NMF，同时能够鼓励相连gene在program space中具有相似loading。
\end{note}

\subsubsection{Latent factor residual dependence NB-NMF}
为了在likelihood层面描述未被$WH$吸收的gene--gene相关性，可以对每个样本引入低维residual latent factor。给定正整数$q<n$，对第$j$个样本令：
\begin{equation*}
	z_j\overset{\mathrm{i.i.d.}}{\sim}\operatorname{N}(0,I_q),\quad z_j\in\mathbb{R}^{q},\quad j=1,2,\dots,p.
\end{equation*}
再令residual loading matrix为：
\begin{equation*}
	B=\begin{pmatrix}
		b_1^{\top} \\
		b_2^{\top} \\
		\vdots \\
		b_n^{\top}
	\end{pmatrix}\in\mathbb{R}^{n\times q},\quad b_i\in\mathbb{R}^{q}.
\end{equation*}
$B$的第$i$行为$b_i^{\top}$，表示第$i$个gene对$q$个residual latent factor的loading。给定$z_j$后，定义条件均值：
\begin{equation*}
	\mu_{ij}(z_j)=(WH)_{ij}\exp(b_i^{\top}z_j),\quad i=1,2,\dots,n,\quad j=1,2,\dots,p,
\end{equation*}
并假设：
\begin{equation*}
	x_{ij}\mid z_j,W,H,B,\phi\overset{\mathrm{ind}}{\sim}
	\operatorname{NB}\left(\mu_{ij}(z_j),\phi_i\right).
\end{equation*}
这里$(WH)_{ij}$是$z_j=0$时的baseline NMF均值；由于$\operatorname{E}[\exp(b_i^{\top}z_j)]=\exp(||b_i||_2^2/2)$，它一般不等于将$z_j$边际化之后的均值。

记标准$q$维正态密度为：
\begin{equation*}
	\varphi_q(z)=(2\pi)^{-q/2}\exp\left(-\frac{1}{2}z^{\top}z\right),\quad z\in\mathbb{R}^{q}.
\end{equation*}
由于不同样本的$z_j$相互独立，观测数据的边际似然为：
\begin{align*}
	L_{\operatorname{latent}}(W,H,\phi,B;X)
	=\prod_{j=1}^{p}\int_{\mathbb{R}^{q}}
	\left\{
	\prod_{i=1}^{n}
	f_{\operatorname{NB}}\left(
		x_{ij};(WH)_{ij}\exp(b_i^{\top}z),\phi_i
	\right)
	\right\}\varphi_q(z)\dif z.
\end{align*}
相应的边际对数似然为：
\begin{align*}
	\ell_{\operatorname{latent}}(W,H,\phi,B;X)
	=\sum_{j=1}^{p}\log\int_{\mathbb{R}^{q}}
	\left\{
	\prod_{i=1}^{n}
	f_{\operatorname{NB}}\left(
		x_{ij};(WH)_{ij}\exp(b_i^{\top}z),\phi_i
	\right)
	\right\}\varphi_q(z)\dif z.
\end{align*}
记：
\begin{equation*}
	\mathcal{C}_{\operatorname{latent}}
	=\left\{(W,H,\phi,B):(W,H,\phi)\in\mathcal{C}_{\operatorname{NB}},\;B\in\mathbb{R}^{n\times q}\right\}.
\end{equation*}
则latent factor residual dependence NB-NMF的极大似然估计问题为：
\begin{equation*}
	(\widehat{W}_{\operatorname{latent}},\widehat{H}_{\operatorname{latent}},
	\widehat{\phi}_{\operatorname{latent}},\widehat{B})
	\in\operatorname*{argmax}_{(W,H,\phi,B)\in\mathcal{C}_{\operatorname{latent}}}
	\ell_{\operatorname{latent}}(W,H,\phi,B;X).
\end{equation*}
该模型继承$W$与$H$的尺度不识别性；又因为$z_j$服从球对称标准正态分布，对任意$q$阶正交矩阵$Q$，$B$与$BQ$给出相同的边际模型。因此，若需要解释单个residual latent factor，还应对$W,H,B$施加适当的识别约束。

\begin{note}
	给定$z_j$后，同一样本内的gene按照逐元素负二项模型条件独立；将共同的$z_j$积分掉之后，这些gene一般不再边际独立。事实上，当$i\neq i'$时，由条件独立性和全协方差公式可得：
	\begin{align*}
		\operatorname{Cov}(x_{ij},x_{i'j})
		={}&(WH)_{ij}(WH)_{i'j}
		\exp\left\{\frac{1}{2}\left(||b_i||_2^2+||b_{i'}||_2^2\right)\right\} \\
		&\times\left\{\exp(b_i^{\top}b_{i'})-1\right\}.
	\end{align*}
	因此，该模型通过共享residual latent factor在likelihood层面诱导residual gene dependence。
\end{note}

\paragraph{变分下界}
边际似然中的$q$维积分通常没有闭式解。对第$j$个样本，取变分分布$Q_j=\operatorname{N}(m_j,S_j)$，其中$m_j\in\mathbb{R}^{q},\;S_j=S_j^{\top}\succ0$，并记其密度为$q_j(z)$。由Jensen不等式可得：
\begin{derivation}
	记$x_{\cdot j}=(x_{1j},x_{2j},\dots,x_{nj})^{\top}$，则：
	\begin{align*}
		\log p(x_{\cdot j})
		&=\log\int_{\mathbb{R}^{q}}
		\frac{p(x_{\cdot j}\mid z)p(z)}{q_j(z)}q_j(z)\dif z \\
		&\geqslant\operatorname{E}_{Q_j}
		\left[\log p(x_{\cdot j}\mid z_j)\right]
		-\operatorname{KL}\left\{Q_j\|\operatorname{N}(0,I_q)\right\} \\
		&=\operatorname{E}_{Q_j}\left[
		\sum_{i=1}^{n}\log f_{\operatorname{NB}}\left(
			x_{ij};(WH)_{ij}\exp(b_i^{\top}z_j),\phi_i
		\right)
		\right]
		-\operatorname{KL}\left\{Q_j\|\operatorname{N}(0,I_q)\right\}.
	\end{align*}
\end{derivation}

记变分参数集合为$\mathcal{V}=\{(m_j,S_j):j=1,2,\dots,p\}$。对所有样本求和得到evidence lower bound：
\begin{align*}
	\mathcal{L}_{\operatorname{ELBO}}(W,H,\phi,B,\mathcal{V})
	={}&\sum_{j=1}^{p}\operatorname{E}_{Q_j}\left[
	\sum_{i=1}^{n}\log f_{\operatorname{NB}}\left(
		x_{ij};(WH)_{ij}\exp(b_i^{\top}z_j),\phi_i
	\right)
	\right] \\
	&-\sum_{j=1}^{p}\operatorname{KL}\left\{Q_j\|\operatorname{N}(0,I_q)\right\}.
\end{align*}
其中Gaussian KL散度具有闭式：
\begin{equation*}
	\operatorname{KL}\left\{Q_j\|\operatorname{N}(0,I_q)\right\}
	=\frac{1}{2}\left\{\operatorname{tr}(S_j)+m_j^{\top}m_j-\log\det(S_j)-q\right\}.
\end{equation*}
相应的variational MLE问题为：
\begin{align*}
	&(\widehat{W},\widehat{H},\widehat{\phi},\widehat{B},\widehat{\mathcal{V}}) \\
	&\quad\in\operatorname*{argmax}_{\substack{
		(W,H,\phi,B)\in\mathcal{C}_{\operatorname{latent}};\\
		m_j\in\mathbb{R}^{q},\;S_j=S_j^{\top}\succ0,\;j=1,2,\dots,p
	}}
	\mathcal{L}_{\operatorname{ELBO}}(W,H,\phi,B,\mathcal{V}).
\end{align*}
这里优化的是边际对数似然的下界，因此得到的是近似MLE。期望项中含有$\log\{1+\phi_i(WH)_{ij}\exp(b_i^{\top}z_j)\}$，一般仍需使用Gaussian quadrature、reparameterization结合Monte Carlo或其它数值方法计算。

\begin{note}
	三个模型形成逐层扩展的关系。原始NB-NMF是最简单的working likelihood：负二项分布描述count noise，非负低秩均值$WH$解释主要共表达；gene graph regularized NB-NMF不改变逐元素似然，而是通过graph Laplacian penalty把外部gene dependence结构注入$W$的估计；latent factor residual dependence NB-NMF则通过共享latent variable的边际化在likelihood层面诱导residual gene--gene dependence。第三种模型的表达能力更强，但边际似然计算更困难，通常需要Laplace approximation、variational approximation或Monte Carlo approximation。
\end{note}

\subsection{优化问题的性质}
\begin{property}\label{prop:NMFOptimization}
	关于Frobenius目标的\gls{NMF}优化问题具有如下性质：
	\begin{enumerate}
		\item 若固定$H$，则$F(W,H)$关于$W$是凸函数；
		\item 若固定$W$，则$F(W,H)$关于$H$是凸函数；
		\item $F(W,H)$关于$(W,H)$一般不是联合凸函数；
		\item 一般不能保证迭代算法收敛到全局最优解。
	\end{enumerate}
\end{property}
\begin{proof}
	(1)固定$H$时，$WH$关于$W$是线性映射，而平方范数是凸函数，所以$F(W,H)$关于$W$为凸函数。\par
	(2)同理可得。\par
	(3)注意到$WH$中含有$w_{ik}h_{kj}$这样的双线性项，所以目标函数关于$(W,H)$不是一般的凸二次函数。取标量情形$f(w,h)=\frac{1}{2}(x-wh)^2$即可看出其Hessian矩阵不总是半正定，因此一般不具有联合凸性。\par
	(4)由(3)可知，局部极小点、鞍点以及初始化敏感性都可能出现，所以通常只能期待得到一个稳定的局部解。
\end{proof}

\subsection{乘法更新算法}
乘法更新是\gls{NMF}中最经典的算法之一。它的优点是形式简单，并且若初始矩阵严格为正，则每一步更新后仍保持非负。

\begin{method}
	在Frobenius目标下，给定初始矩阵$W^{(0)}>0,\;H^{(0)}>0$，可迭代更新：
	\begin{gather*}
		H^{(t+1)}_{kj}=H^{(t)}_{kj}\frac{[(W^{(t)})^{\top}X]_{kj}}{[(W^{(t)})^{\top}W^{(t)}H^{(t)}]_{kj}} \\
		W^{(t+1)}_{ik}=W^{(t)}_{ik}\frac{[X(H^{(t+1)})^{\top}]_{ik}}{[W^{(t)}H^{(t+1)}(H^{(t+1)})^{\top}]_{ik}}
	\end{gather*}
	其中除法按元素进行。为了避免分母为$0$，实际计算中常在分母加入一个很小的正常数。
\end{method}

\begin{derivation}
	由Frobenius目标的梯度：
	\begin{equation*}
		\nabla_HF(W,H)=W^{\top}WH-W^{\top}X
	\end{equation*}
	可将负梯度方向写为“正部分”和“负部分”的差：
	\begin{equation*}
		-\nabla_HF(W,H)=W^{\top}X-W^{\top}WH
	\end{equation*}
	若使用元素依赖的步长：
	\begin{equation*}
		\eta_{kj}=\frac{h_{kj}}{(W^{\top}WH)_{kj}}
	\end{equation*}
	则梯度下降形式$h_{kj}^{+}=h_{kj}-\eta_{kj}(\nabla_HF)_{kj}$给出：
	\begin{align*}
		h_{kj}^{+}
		&=h_{kj}-\frac{h_{kj}}{(W^{\top}WH)_{kj}}\left[(W^{\top}WH)_{kj}-(W^{\top}X)_{kj}\right] \\
		&=h_{kj}\frac{(W^{\top}X)_{kj}}{(W^{\top}WH)_{kj}}
	\end{align*}
	这正是$H$的乘法更新。$W$的更新同理可由$\nabla_WF(W,H)=WHH^{\top}-XH^{\top}$得到。
\end{derivation}

\begin{algorithm}[H]
	\caption{基于Frobenius范数的NMF乘法更新}
	\begin{algorithmic}[1]
		\State \textbf{Input:} 非负矩阵$X\in\mathbb{R}^{m\times n}$，分解秩$r$，最大迭代次数$T$
		\State \textbf{Output:} 非负矩阵$W,H$
		\State 初始化$W^{(0)}>0,\;H^{(0)}>0$
		\For{$t=0,1,\dots,T-1$}
		\State $H^{(t+1)}\gets H^{(t)}\circ\dfrac{(W^{(t)})^{\top}X}{(W^{(t)})^{\top}W^{(t)}H^{(t)}}$
		\State $W^{(t+1)}\gets W^{(t)}\circ\dfrac{X(H^{(t+1)})^{\top}}{W^{(t)}H^{(t+1)}(H^{(t+1)})^{\top}}$
		\State 对$W^{(t+1)}$的列进行归一化，并将尺度吸收到$H^{(t+1)}$
		\If{目标函数下降幅度小于阈值}
		\State 停止迭代
		\EndIf
		\EndFor
	\end{algorithmic}
\end{algorithm}

\begin{method}
	在KL散度目标下，常用乘法更新为：
	\begin{gather*}
		h_{kj}^{+}=h_{kj}\frac{\sum\limits_{i=1}^{m}w_{ik}\dfrac{x_{ij}}{(WH)_{ij}}}{\sum\limits_{i=1}^{m}w_{ik}},\quad
		w_{ik}^{+}=w_{ik}\frac{\sum\limits_{j=1}^{n}h_{kj}\dfrac{x_{ij}}{(WH)_{ij}}}{\sum\limits_{j=1}^{n}h_{kj}}
	\end{gather*}
	当$x_{ij}=0$时，分子中对应项视为$0$。该更新同样保持非负性，并适合稀疏计数矩阵。
\end{method}

\subsection{交替最小二乘与投影梯度}
乘法更新虽然简单，但在接近驻点时可能收敛较慢。另一类常用方法是交替最小二乘：固定$W$时求解非负最小二乘问题，固定$H$时再求解另一个非负最小二乘问题。

\begin{method}
	Frobenius目标下的交替非负最小二乘可写为：
	\begin{gather*}
		H^{(t+1)}=\operatorname*{argmin}_{H\geqslant0}||X-W^{(t)}H||_F^2 \\
		W^{(t+1)}=\operatorname*{argmin}_{W\geqslant0}||X-W(H^{(t+1)})||_F^2
	\end{gather*}
	每个子问题都是凸的非负最小二乘问题，因此可以使用坐标下降、主动集方法或投影梯度法求解。
\end{method}

\begin{method}
	投影梯度法的基本形式为：
	\begin{gather*}
		H^{(t+1)}=\left[H^{(t)}-\eta_t\nabla_HF(W^{(t)},H^{(t)})\right]_+ \\
		W^{(t+1)}=\left[W^{(t)}-\gamma_t\nabla_WF(W^{(t)},H^{(t+1)})\right]_+
	\end{gather*}
	其中$[A]_+$表示逐元素投影到非负半轴，即$([A]_+)_{ij}=\max\{a_{ij},0\}$。步长$\eta_t,\gamma_t$可通过回溯线搜索或固定上界选取。
\end{method}

\subsection{正则化与约束变体}
实际数据分析中，基础\gls{NMF}常常需要加入额外结构。常见形式为：
\begin{equation*}
	\min_{W,H\geqslant0}\frac{1}{2}||X-WH||_F^2+\lambda_W\mathcal{R}_W(W)+\lambda_H\mathcal{R}_H(H)
\end{equation*}
其中$\mathcal{R}_W,\mathcal{R}_H$用于表达额外偏好。

\begin{note}
	常见正则化包括：
	\begin{enumerate}
		\item 稀疏约束：取$\mathcal{R}_H(H)=\sum\limits_{k,j}|h_{kj}|$或$\mathcal{R}_W(W)=\sum\limits_{i,k}|w_{ik}|$，使得每个样本只由少数因子解释，或每个因子只作用于少数变量；
		\item 平滑约束：当变量或样本具有空间、时间顺序时，可惩罚相邻元素差异，使基向量或系数随位置平滑变化；
		\item 正交约束：加入$W^{\top}W\approx I_r$或$HH^{\top}\approx I_r$，使不同因子之间更容易区分；
		\item 图正则化：若样本之间有相似图，可惩罚$\sum\limits_{i,j}a_{ij}||h_i-h_j||^2$，使相似样本具有相近的低维表示。
	\end{enumerate}
	这些约束会改变估计结果的含义，因此不能只按重构误差选择模型，还需要结合解释目标。
\end{note}

\subsection{分解秩的选择}
分解秩$r$控制模型复杂度。$r$太小会把多个结构混合在一起，导致欠拟合；$r$太大则会把噪声或偶然模式解释为因子，导致过拟合和解释不稳定。

\begin{method}
	常见的$r$选择方法包括：
	\begin{enumerate}
		\item 重构误差曲线：计算不同$r$下的$||X-\hat{W}_r\hat{H}_r||_F^2$或$D(X\|\hat{W}_r\hat{H}_r)$，使用类似肘部法则的思路选择误差下降明显变慢的位置；
		\item 交叉验证：随机遮盖部分矩阵元素，用其余元素拟合\gls{NMF}，再比较被遮盖元素上的预测误差；
		\item 稳定性分析：对数据重抽样或对初始值重复运行，若某个$r$下得到的因子结构稳定，则说明该秩更可信；
		\item 领域解释：若问题本身要求固定数量的主题、亚型、成分或部件，则$r$也可由研究目标确定。
	\end{enumerate}
\end{method}

\subsection{预处理与解释}
\begin{note}
	使用\gls{NMF}时需要特别注意如下问题：
	\begin{enumerate}
		\item 非负性是模型的基本条件。若原始数据有负值，应先判断负值是否具有实际意义；若负值是中心化造成的，则通常不应直接套用\gls{NMF}；
		\item 不建议像PCA那样对数据做均值中心化，因为中心化会破坏非负性。常见预处理是尺度归一化、总量归一化、对数变换后平移到非负区域等；
		\item 因子解释依赖尺度。若变量量纲差异很大，大尺度变量会主导重构误差，因此需要根据研究目的决定是否标准化或加权；
		\item 零值具有特殊含义。乘法更新中若某个元素被更新为$0$，之后通常会保持为$0$，所以初始化时常取严格正的小随机数；
		\item \gls{NMF}结果对初始值敏感，实际分析应使用多组初始值重复运行，并选择目标函数较小且解释稳定的解。
	\end{enumerate}
\end{note}

\subsection{应用解释}
\begin{example}
	\textbf{文本主题分析：}令$X$为词项--文档矩阵，$x_{ij}$表示第$i$个词在第$j$篇文档中的出现次数或TF-IDF权重。若$X\approx WH$，则$W$的第$k$列可解释为第$k$个主题下各词的权重，$H$的第$k$行可解释为第$k$个主题在各文档中的强度。
\end{example}

\begin{example}
	\textbf{图像部件表示：}令$X$的每一列是一张灰度图像拉直后的像素向量。若$X\approx WH$，则$W$的列向量可还原成若干基图像，$H$给出每张图像由这些基图像组合的系数。由于系数非负，图像更倾向于被解释为局部部件的叠加。
\end{example}

\begin{example}
	\textbf{生物信息学：}令$X$为基因--样本表达矩阵。$W$的列可视为基因表达模块，$H$的列可视为每个样本在不同模块上的活跃程度。若某些样本在同一模块上系数较大，可能提示这些样本具有相近的生物机制或亚型结构。
\end{example}

\subsection{与其它降维方法的比较}
\begin{table}[H]
	\centering
	\caption{NMF与常见降维方法的比较}
	\begin{tabular}{c|c|c|c}
		\hline
		方法 & 数据约束 & 成分约束 & 主要解释 \\ \hline
		PCA & 一般实值矩阵 & 正交方向 & 最大方差方向 \\
		因子分析 & 一般实值随机向量 & 潜在因子解释协方差 & 公共因子结构 \\
		\gls{NMF} & 非负矩阵 & 非负基与非负系数 & 部件加性组合 \\ \hline
	\end{tabular}
\end{table}

\begin{note}
	\gls{NMF}并不是PCA的简单替代品。PCA有明确的最优低秩逼近理论和正交结构，计算稳定且全局解可由特征值分解或SVD得到；\gls{NMF}牺牲了这类封闭形式解，换来非负约束下更强的可解释性。因此，二者的选择应由数据类型和解释目标决定，而不是只比较重构误差。
\end{note}
